\chapter{Architectural and Organizational Evolution}
\label{chap:architecture}
\chaptermark{Architectural Evolution}

%----------------------------------------------------------------------------------------
% CHAPTER INTRODUCTION
%----------------------------------------------------------------------------------------

This chapter describes the architectural and organizational framework designed to resolve the gap between operational data generation and analytical use. 
Drawing on decentralized  data concepts, semantic modeling, and enterprise knowledge representation, the chapter examines the shift from pipeline-centric 
data architectures toward domain-oriented data products. It then analyzes how corporate ontologies, centralized semantic layers, and knowledge graphs provide 
a reliable foundation for business analytics and AI agents. 
Finally, it outlines the governance processes and organizational changes needed to support this system, focusing on metadata management, the role of Data \& AI Champions, 
and the reorganization of the corporate Data Office.

%----------------------------------------------------------------------------------------
% SECTION 4.1
%----------------------------------------------------------------------------------------

\section[The Data Product Paradigm]{The Data Product Paradigm: Treating Data as a Governed Asset}
\label{sec:data_product_paradigm}

In traditional enterprise architectures, data is often treated as a simple byproduct of daily operations---transactional records generated to support routine tasks and then 
loaded into centralized data warehouses or data lakes. In this setup, downstream analytical teams are left with the task of cleaning, reconciling, and interpreting information 
long after it was created. To address the misalignment between data generation and business consumption, this thesis adopts the Data Product paradigm, a core concept of the 
decentralized Data Mesh framework introduced by Dehghani \cite{dehghani2022data}. In this approach, data is no longer handled as passive operational output; 
instead, it is treated as an independent asset owned by business domains, managed with clear lifecycles, defined quality standards, and consumer-focused design. 
A Data Product includes not only the data itself, but also the transformation code, metadata, access rules, and automated tests needed to make it discoverable, accessible, 
and usable across the enterprise \cite{dehghani2022data}.

In practice, this approach relies on establishing formal Data Contracts and explicit Service Level Agreements (SLAs) between business domains. 
Historically, enterprise IT systems reflect a divide between transactional platforms---managed by the Chief Information Officer (CIO) and focused on operational speed, 
system reliability, and transactional integrity---and analytical platforms, managed by the Chief Data Officer (CDO) to support reporting, forecasting, and machine learning. 
A Data Contract acts as a clear technical agreement designed to manage this difference. 
Instead of allowing data issues to travel downstream into analytical models, the contract defines schema rules, data types, quality limits, and change-management rules 
directly at the source. In this way, data quality shifts from an expensive downstream cleaning task to an architectural requirement enforced from the beginning.

Across an integrated enterprise, this product-oriented model creates clear domain ownership and operational responsibility. In conventional pipelines, 
upstream schema updates or changes to business logic often happen without notice, causing downstream models and reports to fail or report incorrect numbers. 
By packaging data assets as governed Data Products backed by contracts, the enterprise removes this vulnerability. Data producers become directly responsible for schema stability, 
clear definitions, and reliable delivery. Downstream analytical users can then rely on certified, version-controlled interfaces that decouple operational systems from strategic 
decision-support tools, reducing maintenance overhead and increasing trust in shared corporate data.

%----------------------------------------------------------------------------------------
% SECTION 4.2
%----------------------------------------------------------------------------------------

\section[Conceptual Alignment]{Conceptual Alignment: Business Glossary and Ontological Modeling}
\label{sec:business_glossary_ontology}

Before an organization can design stable data systems or automate analytical workflows, it must establish clear common definitions across its departments. 
In large companies, different departments naturally develop their own terms, leading to confusion and misunderstandings. Different business units often use the exact 
same word to describe different things, or use different words for the exact same object. 
According to the governance standards in the \textit{DAMA Guide to the Data Management Body of Knowledge} \cite{dama2024dmbok}, setting up a centralized Business Glossary 
is the standard first step to solve this problem. The Business Glossary works as an enterprise-wide reference that collects approved definitions, clarifies terminology, 
and assigns explicit data owners across business functions \cite{dama2024dmbok}.

However, while a Business Glossary provides clear documentation for human users, a static text dictionary cannot be executed programmatically or used to guide autonomous 
artificial intelligence agents. To enable automated analysis and reliable queries, the enterprise model needs to evolve from a descriptive glossary into a formal, 
machine-readable ontology. 
In information systems research, an ontology is defined as ``a formal, explicit specification of a shared conceptualization'' \cite{gruber1993translation, noy2001ontology}, 
providing a structured framework that machines can process. Unlike a simple glossary, an ontology defines not just business entities and their attributes, 
but also the mathematical, hierarchical, and logical rules connecting them, including class inheritance, functional dependencies, and cardinality constraints.

As literature on data systems shows, database structures and classification rules are never neutral; they reflect design choices and operational limits that determine 
what information is recorded and what is omitted \cite{bowker1999sorting, kitchin2014data}. 
This explains why formal ontological modeling is particularly useful for an integrated agri-food company like Amadori, where standard relational databases and ERP tables 
struggle to represent the complex, time-dependent nature of biological supply chains.
Within Amadori, critical business knowledge has often remained spread across separate operational tools, local spreadsheets, and implicit employee expertise. 
Introducing an enterprise ontology helps the company organize these complex relationships into a unified, computable model. 

By organizing these domain rules into a computable structure, the proposed ontology helps bridge the gap between daily operations and corporate analytics. 
Instead of requiring human planners or AI tools to navigate complex relational tables with obscure column names, an ontology provides a verified semantic baseline. 
Connecting analytical tools and Large Language Model (LLM) agents to this formal structure ensures that automated queries keep the right business context, 
follow operational rules, and avoid the misunderstandings that often occur in complex industrial environments \cite{bowker1999sorting, hilger2026bridging}.

%----------------------------------------------------------------------------------------
% SECTION 4.3
%----------------------------------------------------------------------------------------

\section[The Semantic Layer]{The Semantic Layer: Metric Governance and Decoupled Consumption}
\label{sec:semantic_layer}

While the Business Glossary and corporate ontology establish the conceptual model of enterprise knowledge, the Semantic Layer is the technical component that translates these 
definitions into live, executable queries. It is designed to prevent conflicting definitions and metric discrepancies across different tools. 
In conventional analytical systems, business logic---such as the formulas for promotional uplifts, customer groupings, or baseline sales allocations---is usually scattered 
across individual SQL queries, data transformations, and local reporting dashboards. The Semantic Layer resolves this issue by implementing the abstract definitions 
of the ontology as a centralized, standardized layer placed between physical databases and downstream user interfaces. 
As shown by Hilger, Tesfaye, and Wahl \cite{hilger2026bridging}, this layer acts as the practical link that turns conceptual models into usable analytical tools.

The core principle of the Semantic Layer is the complete separation of raw data storage from business calculation rules. While cloud data platforms and operational systems 
handle storage optimization, data ingestion, and database performance under the IT team's care, the responsibility for defining and validating business metrics 
belongs to the semantic model. By mapping physical database tables and joins into standardized, user-facing business entities, the Semantic Layer creates an enterprise 
single source of truth that shields downstream tools from changes in underlying databases. As Potturi \cite{potturi2025demystifying} notes, 
this architecture establishes a clear data contract between engineering teams and business users, balancing central governance with flexible self-service analysis. 
In this architecture, a business metric is defined only once in the central repository, ensuring that every user or service---whether an operational dashboard, 
a management report, or an analytical model---runs the exact same calculation.

For an integrated manufacturing company like Amadori, a centralized Semantic Layer provides an effective way to resolve discrepancies between high-level sales plans 
and detailed plant operations. In typical operational setups, aligning commercial goals with production limits often requires manual calculations in offline spreadsheets. 
Implementing a Semantic Layer lets the enterprise bring these business rules together into centralized, version-controlled definitions. This shared logic guarantees that 
simulations and operational checks always use identical numbers, regardless of which tool or department runs the query.

Furthermore, the Semantic Layer improves how users and automated systems interact with data by translating business terms into database queries. By hiding complex SQL 
syntax behind clear business metrics, it makes data exploration practical for non-technical domain users. 
As highlighted by Gartner (Richardson et al., 2021) \cite{gartner2021magic}, business intelligence is shifting away from static reports toward interactive tools 
where automated systems and natural language interfaces assist decision-makers. 
Connecting Natural Language Querying (NLQ) and Large Language Model agents to a certified Semantic Layer allows planners to explore connected data domains using normal 
business language \cite{potturi2025demystifying}. Because the AI agent works with verified metric definitions rather than writing raw SQL over complex tables, 
the Semantic Layer acts as a safeguard that prevents calculation errors and hallucinated results, supporting reliable decisions \cite{hilger2026bridging}.

%----------------------------------------------------------------------------------------
% SECTION 4.4
%----------------------------------------------------------------------------------------

\section[Enterprise Knowledge Graphs]{Enterprise Knowledge Graphs: From Abstract Models to Actionable Networks}
\label{sec:knowledge_graphs}

While an ontology defines the general structure and logical rules of a business domain, an enterprise Knowledge Graph is the practical network populated with actual 
enterprise records \cite{hogan2021knowledge}. In semantic modeling, the ontology defines the overall schema, while the Knowledge Graph fills this schema with real 
operational data and connections. In conventional relational databases, data is split across normalized tables, requiring complex and rigid table joins to track multi-step 
business activities. 
A Knowledge Graph addresses this challenge by organizing company data as a network of distinct entities (nodes) linked by explicit relationships (edges), 
making it straightforward to navigate connections, run semantic queries, and trace dependencies \cite{hogan2021knowledge}.

In an integrated supply chain like Amadori's, operational decisions involve interconnected factors that link customer demand, production limits, and raw material availability. 
In traditional systems, evaluating how a change in commercial demand affects production lines requires combining data from different transactional tools, local databases, 
and manual planning tables. Introducing an enterprise Knowledge Graph allows the company to turn its conceptual ontology into an interconnected, queryable network. 
In this structure, commercial demand, daily production schedules, and physical plant assets are linked by explicit relationships. 
This graph model captures the end-to-end operational flow of the business, enabling complete traceability without trying to force complex operational paths into flat database tables.

This graph-based approach also provides a reliable structure for modern artificial intelligence interfaces. 
In typical generative AI setups, system knowledge is added through standard Retrieval-Augmented Generation (RAG), which cuts text documents into small chunks and finds 
information using simple text similarity \cite{lewis2020retrieval}. While effective for answering simple questions from documents, standard RAG struggles 
with multi-step reasoning or following chains of dependencies across different business systems \cite{edge2024local}. 
Because text searches do not understand structured relationships, large language models can easily invent links or jump to false conclusions when connecting separate enterprise domains.

Graph-Augmented Retrieval (GraphRAG) solves this problem by connecting vector searches to the explicit paths of the Knowledge Graph \cite{edge2024local, hogan2021knowledge}. 
Instead of retrieving disconnected text snippets, GraphRAG follows verified relationships in the graph to find connected subgraphs directly related to the user's question. 
Giving the generative model this explicit, verified network of operational facts gives it the solid context needed for complex reasoning \cite{hilger2026bridging}. 
This allows business planners to ask questions about the supply chain in natural language, confident that the AI agent follows real enterprise connections rather than 
guessing relationships based on statistical patterns.

%----------------------------------------------------------------------------------------
% SECTION 4.5
%----------------------------------------------------------------------------------------

\section[Active Metadata Management]{Active Metadata Management: Lineage, Traceability, and Quality Assurance}
\label{sec:metadata_management}

In an architecture using decentralized Data Products, formal ontologies, and semantic layers, metadata is no longer just static documentation; 
it becomes an active operational tool for governance and system reliability. Enterprise architectures are always evolving: underlying database tables, ERP settings, 
and calculation formulas change regularly to match new business procedures and operational needs. 
Without a centralized, active metadata system, an organization builds up technical debt. When that happens, analytical applications and forecasting pipelines 
often fail or report inaccurate results because upstream operational systems changed a column name, data type, or business category without informing downstream teams. 
Under the standards established in the \textit{DAMA Guide to the Data Management Body of Knowledge} \cite{dama2024dmbok}, modern data governance requires moving from passive 
documentation files to active metadata management, where metadata is continuously updated, tracked, and used to run data processes automatically.

The need for systematic metadata management stems from the basic challenges of reusing operational records for analytical purposes. 
As Christine Borgman \cite{borgman2015big} points out, data recorded for routine operational tasks is rarely ready for analytical analysis without careful cleaning and context. 
This problem is also examined by Sabina Leonelli \cite{Leonelli2016data} in her work on data journeys, which explains that when data moves across departmental or technical boundaries, 
it loses the background, assumptions, and practical constraints under which it was recorded. 
To avoid wrong analytical conclusions, data moving from operational tools to planning platforms must include clear metadata describing where it came from, 
how often it is updated, what rules it follows, and any exceptions present when it was recorded \cite{borgman2015big, Leonelli2016data}.
Under the proposed framework, an automated data catalog regularly inspects company pipelines, building an end-to-end data lineage that traces physical sources through 
the Semantic Layer down to dashboards and AI interfaces. When a schema change or calculation update occurs in an upstream system, the metadata platform automatically 
checks for downstream impacts, alerting data owners and preventing unexpected breaks in planning applications.

This automated lineage tracking ensures that data quality is monitored at the source rather than fixed later when problems appear in final reports. 
Ultimately, by maintaining a reliable record of where data originates, how it changes, and how it is classified, active metadata management builds trust in enterprise data. 
It ensures that both human planners and conversational AI systems can read, understand, and use corporate information with complete visibility \cite{borgman2015big, dama2024dmbok}.

%----------------------------------------------------------------------------------------
% SECTION 4.6
%----------------------------------------------------------------------------------------

\section[Organizational Transformation]{Organizational Transformation: The Modern Data Office and Data \& AI Champions}
\label{sec:organizational_transformation}

Introducing decentralized Data Products, centralized semantic models, and conversational AI tools cannot succeed through technical tools alone. 
As enterprise research on analytics shows, technical changes must be supported by practical organizational changes and clear roles if an organization wants to get lasting value 
from its data \cite{davenport2017competing}. 
In traditional organizations, a clear gap often separates technical engineering teams from business decision-makers. 
Data engineers understand database design, cloud systems, and pipeline architecture, but they rarely know the day-to-day business routines and commercial trade-offs 
that give meaning to corporate data. On the other hand, business planners know their daily operations thoroughly, but lack the technical background needed to turn their 
analytical questions into precise software specifications.

To bridge this operational gap, the proposed architecture introduces the role of Data \& AI Champions. These individuals act as practical coordinators working directly 
between business departments and data engineering teams. Following the organizational knowledge concepts described by Nonaka and Takeuchi \cite{nonaka1995knowledge}, 
the main role of these champions is to take the practical, hands-on experience of operational staff---the tacit knowledge---and turn it into clear, documented rules
---the explicit knowledge. Rather than working solely as data engineers or purely as business planners, Data \& AI Champions understand both sides: 
they know commercial and operational goals, and can translate them into formal data contracts, standard metrics, and ontology relationships. 
They also help colleagues adopt conversational AI tools, helping planners move from reading static reports to asking specific analytical questions and evaluating automated 
answers critically.

This operational bridge is supported by a reorganization of the corporate Data Office itself. 
Historically, company data teams worked as a simple support unit inside the IT department under the Chief Information Officer, acting as a help desk responsible 
for writing custom SQL queries and building individual reports on request. 
This older approach inevitably creates an operational bottleneck: whenever business planners need a new metric or an exploratory scenario, 
they have to submit a support ticket to a long IT backlog, delaying decisions.

\begin{table}[t]
\centering
\small
\caption{Separation of Responsibilities Between Technical Systems and Semantic Governance}
\label{tab:governance_responsibilities}
\begin{tabular}{p{3.2cm} p{5.5cm} p{5.5cm}}
\hline
\textbf{Dimension} & \textbf{Technical Custody (CIO Core)} & \textbf{Semantic Governance (CDO Core)} \\ \hline
\textbf{Primary Objective} & Transactional execution, data storage, schema enforcement, and operational performance. & Metric validation, standardized definitions, and strategic decision support. \\ 
\textbf{Core Accountability} & System uptime, infrastructure stability, network security, and database backups. & Consistent definitions, Data Contract compliance, data lineage, and model accuracy. \\ 
\textbf{Governed Layer} & Physical tables, relational databases, system indexes, and transactional logs. & Business glossaries, corporate ontologies, Data Contracts, and standardized metrics. \\ \hline
\end{tabular}
\end{table}

Under the proposed model, the Data Office led by the Chief Data Officer shifts from a service team that simply builds reports to a governance team that verifies and approves standards. 
Instead of creating individual dashboards manually, the reorganized Data Office defines, maintains, and approves the enterprise Business Glossary, the formal ontology, 
and the calculation rules in the Semantic Layer. 
As summarized in \Cref{tab:governance_responsibilities}, hardware management, database performance, and cloud systems stay under the technical care of the CIO, 
while consistent definitions, metric certification, and Data Contract oversight are managed by the CDO. 
By moving from manual report building to managing shared semantic assets, this organizational model removes technical bottlenecks. 
It allows business teams, supported by their champions, to safely explore certified Data Products, while AI agents work with a verified, structured representation of 
enterprise operations \cite{davenport2017competing, dehghani2022data}.